\section{Does Learning Idioms Improve Cultural Understanding?}
\label{sec:experiments}

\subsection{Experiment Setup}
\label{sec:setup}

\paragraph{Training conditions.}
We continue pretraining Qwen3.5-9B-Base \citep{yang2025qwen3} \todo{check that this is the right citation for Qwen3.5} separately for Hindi, Chinese, and Arabic under three conditions and compare with the untouched \textbf{Base} model.
\textbf{\idiomcpt{}} trains on the idiom-annotated corpus of \S\ref{sec:data-corpus}: web documents that contain \dataname{} idioms, each followed by a meaning tag that lists the figurative meanings of the idioms in the document.
\textbf{\random{}} trains on documents sampled at random from the same web sources with no idiom filter and no tags, matched to \idiomcpt{} in number of training tokens. It measures what a model gains from reading more text in the language.
\textbf{\idiomdocs{}} trains on exactly the documents of \idiomcpt{} with the meaning tags removed, for the same number of optimization steps. Comparing it with \random{} isolates the effect of \textit{selecting} idiom-bearing documents, and comparing \idiomcpt{} with it isolates the effect of \textit{explaining} the idioms.
All runs use the same recipe: full-parameter training with LLaMA-Factory \citep{zheng2024llamafactory}, 16K-token packed sequences, a global batch of 2.1M tokens, learning rate $10^{-5}$ with cosine decay, and three epochs, which amounts to 2.1K steps for Hindi (1.4B-token corpus), 11.2K for Chinese (7.8B), and 1.6K for Arabic (1.1B). Appendix~\ref{app:tokens} gives the token budgets of all runs.

\paragraph{Benchmarks.}
We evaluate on 17 multiple-choice benchmarks, 16 of them in five groups that are increasingly distant from the training signal.
(1) \textit{Idiom meaning}: choosing the meaning of an idiom, with Kinayat-Meaning for Arabic and the connotation task of Chengyu-Bench for Chinese \citep{kinayat2025,chengyubench2025}. These idioms and their meanings appear in the meaning tags, so this group measures whether the knowledge is absorbed.
(2) \textit{Figurative inference}: AR-Figurative, a 314-item set that we assemble from the figurative-language subsets of Alyah and DziriEval and in which about 99\% of the expressions are \textit{absent} from \dataname{}, and the Hindi portion of MABL \citep{kabra2023multi}. This group measures generalization to unseen figurative language.
(3) \textit{Idiom cloze}: selecting the idiom that fits a context, with Kinayat-Cloze and ChID \citep{zheng2019chid}.
(4) \textit{Culture}: ArabCulture \citep{sadallah2025arabculture}, ArabicCulturalQA \citep{arabicculturalqa2025}, Alyah \citep[Emirati culture;][]{alyah2025}, DziriEval \citep[Algerian culture;][]{dzirieval2025}, Global-PIQA \citep{chang2025global} in Arabic and Hindi, and the classical poetry matching task CCPM \citep{li2021ccpm}. This group is the test of our hypothesis.
(5) \textit{Regional knowledge}: ArabicMMLU \citep{koto2024arabicmmlu}, MILU \citep{verma2025milu}, and CMMLU \citep{li2023cmmlu}.
The seventeenth benchmark is a culture-agnostic control, the parallel (translated) split of Global-PIQA in Arabic, on which no trained model differs significantly from the base model.
All models are scored with the same prompts by the log-likelihood of each option. Because all conditions answer the same items, we report item-level paired bootstrap 95\% confidence intervals and McNemar tests for every contrast.

\subsection{Do Models Learn Idioms from Idiom-Centered Pretraining?}
\label{sec:results-idiom}

\input{latex/tables/main_results.tex}

Yes, and they learn more than what is written in the tags. Table~\ref{tab:main} and Figure~\ref{fig:cpt-effects}(a) report all conditions. On idiom meaning, \idiomcpt{} improves over \random{} by 21.5 points on Kinayat-Meaning (23.7 to 45.2) and by 6.9 points on Chengyu-Bench (88.9 to 95.7). These benchmarks overlap with the tags by design, so the more informative result is AR-Figurative: although its expressions are almost entirely outside the knowledge base, \idiomcpt{} gains 5.4 points over \random{}, while \random{} itself falls slightly below the base model. Reading idioms together with their meanings therefore improves the interpretation of figurative expressions the model was never taught. The gain does not extend to idiom cloze. On ChID and Kinayat-Cloze all three trained models are within a point of each other, which suggests that choosing an idiom for a context depends on distributional fit that any in-language text provides, and knowing what an idiom means adds little. Hindi MABL follows the same pattern: a 6.1-point gain over the base model that \random{} matches. \todo{Hindi has no idiom-meaning benchmark in this table, so the ``knowledge is absorbed'' claim rests on ar and zh. Options: build a meaning-selection test from held-out \dataname{} Hindi proverbs, since the only other Hindi idiom measure, LLM-judged English-to-Hindi idiom translation (IdiomCE), improves over the base model (2.06 to 2.48, $p<10^{-10}$) but ties with \random{} (2.48 vs.\ 2.50). In Hindi, no contrast between \idiomcpt{} and \random{} is significant; the old draft attributes this to the weak base model, which is a hypothesis we have not tested.}

\input{latex/figures/fig_cpt_effects.tex}

\subsection{Does Idiom Knowledge Transfer to Cultural Understanding?}
\label{sec:results-culture}

Only in isolated cases. Compared with the base model, \idiomcpt{} looks like a clear cultural improvement: it gains significantly on Alyah ($+4.7$), ArabCulture ($+3.4$), DziriEval ($+2.6$), and CCPM ($+2.2$), as well as on MILU ($+7.6$). The token-matched control changes the conclusion. \random{} obtains nearly the same gains, and the difference between \idiomcpt{} and \random{} is significant on only one of the seven culture benchmarks, Alyah ($+2.9$, 95\% CI $[1.2, 4.6]$). DziriEval ($+1.8$) and Hindi Global-PIQA ($+7.0$ on 100 items) point in the same direction without reaching significance, and the remaining four are within about a point of zero. On regional knowledge the three trained models are indistinguishable. A study that compared only against the base model, as is common practice, would have attributed all of these gains to idioms.

We draw two conclusions. First, most of what current culture benchmarks measure is acquired from generic text in the language, and curating that text around idioms neither helps nor hurts. Second, the places where idioms do help are informative. Alyah and DziriEval test the culture of specific dialect communities (Emirati and Algerian), and both contain a figurative-language subset, the source of our AR-Figurative set. The culture benchmarks on which idioms help are thus the ones in which proverbs and sayings are part of what is being tested. \todo{Important check: recompute Alyah and DziriEval \textit{excluding} their figurative subsets. If the Alyah gain disappears, the honest statement is that idiom training helps zero of seven culture benchmarks once figurative items are removed; if it remains, the second explanation applies: 55\% of the Arabic knowledge base is colloquial regional proverbs, and idiom-bearing documents may be where regional culture is written down (test by measuring the share of dialectal documents in the \idiomcpt{} and \random{} corpora).} \todo{No correction for multiple comparisons is applied across the 17 benchmarks. Under Holm correction over all 17, Kinayat-Meaning, Chengyu-Bench, and Alyah ($p=0.0014$) survive, but AR-Figurative ($p=0.0076$ against a threshold of $0.0036$) does not; decide on the family of tests (e.g. per benchmark group) and report corrected values.}

\subsection{Where Does the Gain Come From?}
\label{sec:results-decomp}

Figure~\ref{fig:cpt-effects}(b) decomposes the gain of \idiomcpt{} over \random{} on the benchmarks where it is largest. Selecting idiom-bearing documents, without any explanation, already accounts for most of the gain on AR-Figurative (3.8 of 5.4 points), Alyah (2.1 of 2.9), and Chengyu-Bench (4.3 of 6.9): idioms in context carry much of their own meaning. The meaning tags add a significant amount only on the two idiom-meaning benchmarks, 13.5 points on Kinayat-Meaning and 2.6 on Chengyu-Bench, where the tag states the answer. Explicit definitions are thus an efficient way to install dictionary knowledge, but the transfer to unseen figurative language and to regional culture comes mainly from the documents.

\subsection{Discussion}
\label{sec:discussion}

\paragraph{Why does idiom knowledge not transfer further?}
Section~\ref{sec:analysis} showed that idioms encode \textit{symbolic and evaluative} knowledge: what a dog, a dragon, or the color red stands for. Culture benchmarks ask about \textit{facts and practices}: dishes, ceremonies, etiquette, history. These are different layers of culture, and a model can gain the first without the second. \todo{This is the central argument of the paper and currently has no experiment behind it. Proposed experiment: build a symbolism probe from the same-entity analysis (e.g. ``In Chinese culture, calling someone a dog implies \ldots'' with options drawn from the en/zh primary meanings; several hundred items can be generated from the 515 entity analyses and verified by annotators), then evaluate Base, \random{}, \idiomdocs{}, and \idiomcpt{} on it. The prediction is a gain of \idiomcpt{} over \random{} on the probe that is absent on fact-based benchmarks. A second test: within existing culture benchmarks, compare accuracy on items that mention high-divergence entities against the rest.}

\paragraph{Does the advantage survive instruction tuning?}
We fine-tune Base and \idiomcpt{} on the same general instruction mixture (60\% target language, 40\% English) and score both with the same log-likelihood protocol. The advantage of the \idiomcpt{} initialization persists on every idiom and culture benchmark in Hindi and Arabic: Kinayat-Meaning 35.1 vs.\ 18.8, AR-Figurative 34.4 vs.\ 30.2, Kinayat-Cloze 54.7 vs.\ 52.0, and Alyah 37.5 vs.\ 35.2 in Arabic; Global-PIQA 69.0 vs.\ 60.0, MILU 58.3 vs.\ 56.7, and MABL 56.3 vs.\ 54.1 in Hindi. In Chinese it persists on the idiom benchmarks (Chengyu-Bench 95.2 vs.\ 90.0, ChID 63.8 vs.\ 62.8), and the two initializations converge on CMMLU and CCPM. The source of the instruction data matters for culture: replacing a machine-translated Arabic mixture (300K examples) with a native Arabic mixture less than a third of its size raises every Arabic idiom and culture benchmark (e.g.\ Kinayat-Meaning 19.1 to 35.1, ArabCulture 46.2 to 49.8, DziriEval 29.8 to 34.7). \todo{(i) No confidence intervals were computed for any instruction-tuned result. (ii) The decisive control, instruction tuning the \random{} checkpoint, was not run; since \random{} matches \idiomcpt{} on most culture benchmarks before instruction tuning, the culture part of this advantage is probably in-language exposure again. (iii) Turn this paragraph into a table (Base-IT, \random{}-IT, \idiomcpt{}-IT) once (i) and (ii) are done.}

\paragraph{What does idiom-centered training cost?}
Monolingual continued pretraining erodes English abilities, most of all code generation, and the erosion grows with the training budget: HumanEval falls from 72.0 to 64.6 for Arabic (3.4B training tokens), 59.2 for Hindi (4.4B), and 53.7 for Chinese (23.4B), while MMLU stays within 1.5 points in all three (Appendix Table~\ref{tab:retention}). This cost is a property of continued pretraining and not of idiom data: \random{} loses the same amount or slightly more (HumanEval 64.0, 56.1, and 53.7). Instruction tuning with the English anchor restores most of it (HumanEval 68.9 for Hindi and 67.7 for Chinese).

\paragraph{Remaining experimental gaps.}
\todo{(1) One base model and one seed; add a second model family or size and at least three seeds for the key contrasts. (2) No English training arm, although English is half of the analysis. (3) ArabCulture accuracy differs between \texttt{FINAL\_REPORT.md} (45.8/49.1/48.7, $n{=}3{,}471$) and \texttt{ci\_report.json} (47.2/50.6/50.3, $n{=}2{,}168$); resolve which split is reported (the old COLM draft uses the former, Table~\ref{tab:main} the latter). (4) ChID base and \idiomcpt{} were run on 3,000 items and the controls on 3,756; paired tests use the shared 3,000. (5) Chengyu-Bench reports only the connotation subtask (540 items); the appropriateness subtask (572 items) is missing. (6) The token match is approximate for Hindi (\random{} processed 4.09B tokens vs.\ 4.40B for \idiomcpt{}, about 7\% fewer), and \random{} corpora contain 1.5--3.3$\times$ more, shorter documents; state this and, if possible, rerun Hindi \random{} to the exact budget. (7) Report every null result in the main text or appendix, including the ones the old draft commented out (ChID, Hindi IdiomCE where \random{} scores 2.50 vs.\ 2.48, and the culture-agnostic control row). (8) Qualitative win/loss examples per idiom benchmark from the stored per-item records would make a useful case-study table here.}
